\clearpage
\section{Özet}
\label{CH:BolumOzet}

Bu projede, işitme ve konuşma engelli bireylerin günlük iletişimde karşılaştığı
sorunları azaltmak amacıyla akıllı telefon kamerasından alınan Türk İşaret Dili
(TİD) hareketlerini tanıyıp yazıya çeviren bir sistem geliştirilmektedir. Sistem,
işareti video olarak kaydedip bir sunucuya göndermek yerine telefonda çıkarılan
iskelet koordinatları üzerinde çalışmaktadır. Böylece internet bağlantısına ihtiyaç
kalmamakta, kamera görüntüsü de cihazdan dışarı çıkmamaktadır.

Veri tarafında AUTSL veri setinin resmî \texttt{train}, \texttt{val} ve
\texttt{test} bölmeleri kullanılmıştır. Bölmeler arasında işaretçi kesişmediği
doğrulanmış, açılamayan 11 video tespit edilerek çıkarım dışında bırakılmıştır.
Her karede iki elin 21'er noktası ve üst gövdeye ait 7 nokta Google MediaPipe
Tasks API ile çıkarılmıştır. Koordinatlar omuz ortasına göre ötelenmiş ve omuz
genişliğine bölünmüş, değişken uzunluktaki diziler 64 kareye örneklenmiştir.

Sınıflandırma aşamasında PyTorch ile iki katmanlı, iki yönlü bir GRU (BiGRU)
temel modeli ve karşılaştırma amacıyla üç katmanlı 1B-CNN + iki katmanlı
Transformer modeli kurulmuştur. Aynı veri ve artırma düzeninde 60 epoch sonunda
BiGRU modeli test kümesinde yüzde 88,24 (doğrulama yüzde 89,23), daha hafif olan
1B-CNN + Transformer modeli ise yüzde 88,83 (doğrulama yüzde 89,97) Top-1
doğruluğa ulaşarak tez önerisindeki yüzde 85 hedefini aşmıştır. Temel modelin
LiteRT int8 sürümü 3,19 MB boyutundadır ve aynı 200 test örneği Python ile
telefon üzerinde çalıştırıldığında tahminlerin yüzde 99,5'i birbiriyle aynı
çıkmıştır.

Mobil tarafta Kotlin/CameraX ile kurulan kamera ve çıkarım hattı Flutter
arayüzüne bağlanmıştır. Telefonda model çıkarımı ortalama 48 ms sürmektedir.
Buna karşılık kamera, MediaPipe ve arayüz dâhil uçtan uca hız yaklaşık 9 FPS'tir
ve hedeflenen 15 FPS'in altında kalmaktadır. Canlı kamera kullanımında tahmin
güvenlerinin kaydedilmiş veriye göre belirgin biçimde düştüğü gözlenmiş, nedeni
henüz kesinleştirilememiştir. Bu iki konu sonuç bölümünde ayrıca tartışılmaktadır.

\clearpage
